\section{Executable certificate protocol}
\subsection{Problem objects}
Integers are JSON integers, not Booleans. A rational is a pair \code{[numerator, denominator]} with a strictly positive denominator. The parser canonicalizes rational values and polynomial terms. Vector lengths must agree with the problem dimension.

A Petri input has exactly the following top-level fields:
\begin{lstlisting}
{
  "kind": "petri",
  "initial": [1, 0, 0],
  "transitions": [
    {"pre": [1, 0, 0], "post": [0, 1, 0]},
    {"pre": [0, 1, 0], "post": [1, 0, 0]},
    {"pre": [0, 0, 0], "post": [0, 0, 1]}
  ],
  "targets": [[0, 2, 0]]
}
\end{lstlisting}
The cubic branching input is:
\begin{lstlisting}
{
  "kind": "pps",
  "polynomials": [[
    {"coefficient": [1, 4], "powers": [0]},
    {"coefficient": [3, 4], "powers": [3]}
  ]]
}
\end{lstlisting}
The \code{subject} field in a receipt is the complete canonical JSON encoding of the normalized problem, not a hash or a pretty-printed summary. The checker receives the problem independently and compares the exact encoding. This binds Python acceptance to the finite mathematical input; a future source bridge must additionally bind that input to the original Lean expression and context.

\subsection{Accepted receipt kinds}
\begin{longtable}{P{0.22\textwidth}P{0.68\textwidth}}
\toprule
Kind & Mathematical payload \\
\midrule
\code{safe} & \code{basis}, \code{target_cover}, and \code{predecessor_cover}; indices certify the inclusions in Theorem~\ref{thm:safety}.\\[4pt]
\code{coverable} & \code{word}, a finite list of transition indices; replayed from the original initial marking.\\[4pt]
\code{enclosure} & \code{ledger}, final \code{lower} and \code{upper}, \code{requested_bits}, and a rechecked \code{target_met}. An optional \code{contraction} contains \code{weight} and \code{rho}.\\[4pt]
\code{extinction_one} & A positive rational \code{weight}; normalization, graph connectivity, mean inequality, and nondegeneracy are recomputed.\\[4pt]
\code{algebraic_lfp} & A nested accepted scalar \code{enclosure}, a nonconstant \code{factor}, and a nonzero \code{cofactor}, each polynomial listed in ascending degree.\\
\bottomrule
\end{longtable}

Ledger steps are one of three forms. A \code{kleene} step supplies \code{next}. A \code{newton} step supplies \code{inverse} and \code{next}. A \code{weighted_newton} step supplies \code{weight}, \code{direction}, and \code{next}. All final endpoints are decoded and compared with the endpoints actually obtained from the ledger.

The \code{unknown} kind is a search outcome and is never accepted as a certificate. Statistics such as expansion counts are not mathematical proof payload. The checker does not trust them, and changing them is not expected to invalidate a correct theorem.

\subsection{Input rejection and resource control}
The parser and checker validate dimensions, signs, indices, and the exact inequalities stated in the report. They are research code, not a hardened public service. Extremely large integers, exponents, matrices, nested JSON, or ledgers can consume excessive memory or time. Duplicate-key rejection, byte limits, aggregate term-degree limits, and process-level resource isolation belong in a production ingestion layer.

The producer's work budgets do not replace parser/checker budgets. An externally supplied receipt might be much larger than anything the local producer would emit. In Lean, totality of a recursive checker does not imply that reducing it is affordable. Admission limits therefore need to be specified independently of mathematical validity.

\section{Checker-oriented proof dependency graph}
The Petri lane depends on just three mathematical facts: componentwise predecessor equivalence, complement invariance under a backward-closed set, and induction on finite executions. Dickson's lemma belongs only to the later search-completeness proof.

The enclosure lane depends on valid polynomial monotonicity, existence of the least fixed point, the prefixed upper lemma, and a ledger induction. The weighted variant additionally uses the directional remainder and finite maximum principle. It does not require formalizing exact Gaussian elimination or proving the supplied matrix is an inverse.

The critical lane uses normalization, the finite product-deficit inequality, graph propagation of a maximizing ratio, and a strictly positive deficit supplied either by a strict mean row or by actual branching. The algebraic decoder adds a scalar factor identity, interval evaluation soundness, and strict monotonicity from the derivative sign.

This graph is an implementation aid: a failed probability-source bridge should not prevent proving the finite weighted-step theorem, and a missing WQO library lemma should not block an individual Petri safety theorem.
